%%%%%%%%%%%%
%
% $Autor: Wings $
% $Datum: 2019-03-05 08:03:15Z $
% $Pfad: SetUp.tex $
% $Version: 4250 $
% !TeX spellcheck = en_GB/de_DE
% !TeX encoding = utf8
% !TeX root = manual 
% !TeX TXS-program:bibliography = txs:///biber
%
%%%%%%%%%%%%

\chapter{Bedienung}

Nach der Inbetriebnahme kann ELMAR für den Schneidvorgang verwendet werden.

\section{Bedienung}

\begin{enumerate}
	\item ELMAR auf eine feste, ebene und nicht brennbare Fläche stellen.
	
	```
	\item Den Fahrbereich der Achsen freiräumen und prüfen, ob sich nichts im Bewegungsbereich befindet.
	
	\item Den Styropor-Rohling in den Arbeitsbereich einsetzen und sicher befestigen.
	
	\item Laptop und ELMAR über das USB-Kabel verbinden.
	
	\item ELMAR einschalten.
	
	\item UGS auf dem Laptop starten.
	
	\item In UGS den passenden COM-Port auswählen.
	
	\item Die Baudrate auf 115200 einstellen.
	
	\item Über \textit{Connect} die Verbindung zur Steuerung herstellen.
	
	\item Die Referenzfahrt in UGS starten und warten, bis alle Achsen ihre Referenzposition erreicht haben.
	
	\item Die gewünschte G-Code-Datei in UGS öffnen.
	
	\item Den angezeigten Werkzeugpfad kontrollieren.
	
	\item Die Startposition passend zum Styropor-Rohling anfahren.
	
	\item Den Nullpunkt in UGS setzen, zum Beispiel über \textit{Zero All}.
	
	\item Den Heizdraht einschalten und kurz aufheizen lassen.
	
	\item Prüfen, ob sich niemand im Bewegungsbereich der Maschine befindet.
	
	\item Den Schneidvorgang in UGS starten.
	
	\item ELMAR während des Schneidvorgangs beobachten.
	
	\item Nach dem Schneiden warten, bis keine Bewegung mehr stattfindet.
	
	\item Den Heizdraht abkühlen lassen.
	
	\item Das geschnittene Werkstück aus dem Arbeitsbereich entnehmen.
	
	\item ELMAR ausschalten.
	
	\item Die USB-Verbindung trennen.
	
	\item Materialreste entfernen und den Arbeitsplatz kontrollieren.
	```
	
\end{enumerate}
